\section{Estado del arte}
\label{sec:estado_arte}

Cita aquí tus fuentes usando el comando \texttt{\textbackslash cite}.

La clasificación de textos financieros ha evolucionado desde sistemas basados en reglas hacia modelos de aprendizaje automático capaces de gestionar la ambigüedad del lenguaje natural.

\begin{itemize}
    \item \textbf{Técnicas de vectorización:} Para transformar texto en datos numéricos, destaca el uso de \textit{TF-IDF}, que asigna pesos a los términos según su relevancia en el corpus \cite{kamath2018}. Por otro lado, los \textit{Word Embeddings} permiten capturar relaciones semánticas más complejas al representar palabras en espacios vectoriales continuos \cite{minaee2021}.
    \item \textbf{Modelos de clasificación:} El uso de algoritmos clásicos como \textit{SVM} y \textit{Random Forest} sigue siendo fundamental en entornos que requieren eficiencia computacional y robustez en textos cortos \cite{kolluri2019}. Aunque los \textit{Transformers} (como BERT o RoBERTa) han marcado el estado del arte en precisión, su alta demanda de memoria y capacidad de cómputo plantea desafíos para su despliegue en aplicaciones más ligeras \cite{minaee2021}.
    \item \textbf{Aprendizaje incremental (Online Learning):} En el sector financiero, es crucial que los modelos se adapten a nuevos patrones de gasto sin reentrenamientos completos. El uso de algoritmos optimizados mediante Descenso de Gradiente Estocástico (SGD) permite el aprendizaje incremental, una técnica necesaria para la escalabilidad y el mantenimiento de sistemas en producción \cite{kolluri2019}.
\end{itemize}